\documentclass[11pt]{article}
\usepackage[margin=1in]{geometry}
\usepackage{booktabs}
\title{SNc $\to$ VTA Module Preservation --- Interpretation}
\author{GSE233866, mouse midbrain DA neurons, healthy baseline}
\date{}

\begin{document}
\maketitle

\subsection*{Why this test, not just the two networks side by side}

Independent SNc and VTA networks disagree about which module CACNA1D belongs
to (turquoise vs. brown) and disagree sharply on module count (9 vs. 39) ---
on its own, ambiguous evidence for whether the two regions are structurally
different or whether that is simply an artifact of unequal sample sizes
(2{,}226 vs. 832 cells; see \texttt{../vta\_network/INTERPRETATION.tex} for
why VTA's finer fragmentation is likely such an artifact).
\texttt{ModulePreservation} resolves this by testing a sharper, formal
question: holding SNc's modules fixed as the reference, does that same gene
community reproduce statistically in VTA's data, using 200 permutations to
build a null distribution rather than relying on a raw overlap count.

\subsection*{Result}

The reference (\textbf{turquoise}, CACNA1D's SNc module) reaches
$Z_{\text{summary}} = 20.4$, comfortably above the $Z>10$ strong-preservation
threshold, ranking 3rd of 8 real modules --- behind only brown (53.1) and
green (47.2), and well ahead of yellow (14.9), blue (14.6), red (10.9),
black (7.9), and pink (7.6). This changes the interpretation of CACNA1D's
module-identity switch: rather than evidence of a region-specific network
dissolving, it more likely reflects CACNA1D's own connectivity shifting
\emph{within} a gene neighborhood that remains statistically intact across
both regions. This formal result is independently corroborated by a much
simpler, non-statistical check: 1{,}009 of turquoise's 1{,}352 genes reappear
verbatim in a completely separate SNc network built from a different animal
cohort (\texttt{../../lesion\_model/snc\_intact\_network/}), a 75\% raw
overlap consistent with turquoise being a robust, non-arbitrary gene
community rather than a clustering artifact.

\subsection*{Reading the \texttt{.qual} vs.\ \texttt{.pres} columns}

\texttt{preservation.csv} carries two families of Z-statistics per module,
distinguished by column suffix. \texttt{.qual} columns test whether a
module is well-defined \emph{within its own reference data} --- a
self-consistency check, not a cross-dataset one. \texttt{.pres} columns
test whether the module's structure carries over into the \emph{query}
(VTA) dataset --- this is the one used above and the one that answers the
actual question this test was built for. The two can diverge: a module can
be a clean, well-separated cluster in SNc alone (\texttt{Zsummary.qual}) yet
still fail to reproduce in VTA (\texttt{Zsummary.pres}) if the underlying
biology genuinely differs between regions. That turquoise scores well on
\texttt{Zsummary.pres} specifically is the stronger claim of the two.

\subsection*{Caveats}

The permutation test ran to completion in 12 minutes at 3.88\,GB peak
memory --- comfortably under the resource ceiling that blocked the
equivalent preservation test on the human dataset elsewhere in this
project, so this result did not require any of the memory-saving
compromises that test needed. \texttt{grey} and \texttt{gold} rows in the
raw CSV are WGCNA-internal bins (unassigned genes and a random-gene
statistical control, respectively) and are excluded from the module count
and ranking above. \texttt{moduleSize.obs} in the same file is capped at
1{,}000 genes for every module regardless of true size --- a computational
ceiling on the permutation procedure, not a claim that turquoise (actually
1{,}352 genes) shrank. One sample-size caveat worth carrying forward: VTA's
smallest contributing animal (\texttt{s450}) supplied exactly 102 cells
against a 100-cell minimum threshold --- not disqualifying, but the
preservation result is not resting on a large margin at the low end, and a
slightly stricter QC threshold could have excluded that animal entirely.

\end{document}
